\documentclass[11pt,a4paper]{article}

\usepackage[T1]{fontenc}
\usepackage[utf8]{inputenc}
\usepackage[a4paper,margin=27mm]{geometry}
\usepackage{amsmath,amssymb,amsthm,mathtools}
\usepackage{newpxtext,newpxmath}
\usepackage{microtype}
\usepackage{booktabs,tabularx,array}
\usepackage{enumitem}
\usepackage{xcolor}
\usepackage{titlesec}
\usepackage[authoryear,round]{natbib}
\usepackage[colorlinks=true,linkcolor=blue!45!black,
  citecolor=red!45!black,urlcolor=blue!45!black]{hyperref}

\definecolor{otblue}{HTML}{1F4E79}
\definecolor{otgray}{HTML}{F1F3F5}
\definecolor{otline}{HTML}{AAB5BF}
\titleformat{\section}{\normalfont\Large\bfseries\color{otblue}}
  {\thesection}{.7em}{}
\titleformat{\subsection}{\normalfont\large\bfseries\color{otblue!80!black}}
  {\thesubsection}{.65em}{}
\titleformat{\paragraph}[runin]
  {\normalfont\bfseries\color{otblue!75!black}}{}{0pt}{}
\titlespacing*{\section}{0pt}{2.7ex plus .7ex minus .2ex}{1.1ex}
\titlespacing*{\subsection}{0pt}{2ex plus .5ex minus .2ex}{.75ex}
\titlespacing*{\paragraph}{0pt}{1.3ex plus .3ex minus .1ex}{.55em}

\newtheorem{theorem}{Theorem}[section]
\newtheorem{proposition}[theorem]{Proposition}
\newtheorem{corollary}[theorem]{Corollary}
\newtheorem{lemma}[theorem]{Lemma}
\theoremstyle{definition}
\newtheorem{definition}[theorem]{Definition}
\newtheorem{example}[theorem]{Example}
\theoremstyle{remark}
\newtheorem{remark}[theorem]{Remark}

\newcommand{\R}{\mathbb R}
\newcommand{\Pcal}{\mathcal P}
\newcommand{\Ccal}{\mathcal C}
\newcommand{\Scal}{\mathcal S}
\newcommand{\Cfrak}{\mathfrak C}
\newcommand{\Id}{\operatorname{Id}}
\newcommand{\Lip}{\operatorname{Lip}}
\newcommand{\KL}{\operatorname{KL}}
\newcommand{\BW}{\operatorname{BW}}
\newcommand{\GW}{\operatorname{GW}}
\newcommand{\FR}{\operatorname{FR}}
\newcommand{\argmin}{\operatorname*{arg\,min}}
\newcommand{\dd}{\,\mathrm d}
\newcommand{\eps}{\varepsilon}
\newcommand{\norm}[1]{\left\lVert #1\right\rVert}
\newcommand{\abs}[1]{\left\lvert #1\right\rvert}
\newcommand{\ip}[2]{\left\langle #1,#2\right\rangle}
\newcommand{\cl}{\operatorname{cl}}
\newcommand{\dom}{\operatorname{dom}}
\newcommand{\tr}{\operatorname{tr}}

\hypersetup{
  pdftitle={The Structure of c-Concave Functions},
  pdfauthor={Gabriel Peyre}
}

\title{The Structure of $c$-Concave Functions\\
  \large Envelopes, ordinary convexity, cones, and cross-curvature}
\author{Gabriel Peyr\'e}
\date{July 2026}

\begin{document}
\maketitle

\begin{abstract}
For an arbitrary cost, $c$-concave functions form a lower-envelope class: they
are closed under pointwise infima and addition of constants, hence form a
min-plus, or tropical, cone. They need not form an ordinary convex set. This
note identifies the additional structures that produce ordinary convexity.
For the bilinear cost $c(x,y)=-\langle x,y\rangle$, one recovers the cone of
closed concave functions. For $c(x,y)=\lvert x-y\rvert^2/(2\tau)$, one obtains
the convex, but non-conic, set of $\tau$-semiconcave functions. For the metric
cost $c=Ld$, one obtains the closed unit ball $\operatorname{Lip}(u)\leq L$,
again convex but not a cone. For smooth twisted costs, the theorem of
Figalli--Kim--McCann characterizes convexity of the $c$-concave class by
nonnegative cross-curvature, a strengthening of the weak Ma--Trudinger--Wang
condition. This explains the nontrivial convexity result for squared geodesic
distance on round spheres and products of round spheres. We also summarize
Bregman, quotient, Wasserstein-lifted, matrix, information-geometric and
Gromov--Wasserstein examples, and distinguish ordinary convexity from the
universal tropical structure and from quantitative strong $c$-concavity.
\end{abstract}

\section{Conventions and the universal structure}

The answer depends critically on the sign convention, the domains, and whether
``convex'' refers to ordinary linear combinations or to min-plus combinations.
We use throughout the convention of the OT4ML book.

\subsection{Transform, fixed points, and supports}

Let $X$ and $Y$ be sets and let
$c:X\times Y\to\R\cup\{+\infty\}$. For clarity, most statements below are
written for finite-valued functions and costs; the same identities hold for
extended-real-valued functions whenever undefined combinations of infinities
are excluded.

\begin{definition}[$c$-transform and $c$-concavity]\label{def:c-transform}
For $f:X\to\R\cup\{-\infty\}$ and $u:Y\to\R\cup\{-\infty\}$, define
\[
  f^c(y):=\inf_{x\in X}\{c(x,y)-f(x)\},
  \qquad
  u^{\bar c}(x):=\inf_{y\in Y}\{c(x,y)-u(y)\}.
\]
A function $u$ on $Y$ is $c$-concave if $u=f^c$ for some $f$ on $X$. We denote
the resulting class by $\Cfrak_c(Y)$. The $c$-subdifferential is
\[
  \partial^c u(y)
  :=\{x\in X:u(y)+u^{\bar c}(x)=c(x,y)\}.
\]
Its elements index the cost sections that support the lower envelope at $y$.
\end{definition}

Thus every $c$-concave function is an infimum of shifted cost sections,
\[
  u(y)=\inf_{x\in X}\{c(x,y)-f(x)\}.
\]
These sections play the role of affine supports in ordinary concave analysis.

\begin{proposition}[Closure by double transformation]\label{prop:closure}
The transforms reverse order. Moreover,
\[
  f\leq f^{c\bar c},
  \qquad
  f^{c\bar c c}=f^c.
\]
Consequently, $u$ is $c$-concave if and only if
$u=u^{\bar c c}$. In that case $u^{\bar c c}$ is the smallest $c$-concave
majorant of an arbitrary function $u$.
\end{proposition}

\begin{proof}
If $f_0\leq f_1$, then $c-f_0\geq c-f_1$, hence
$f_0^c\geq f_1^c$. Also, $f^c(y)\leq c(x,y)-f(x)$ for every $(x,y)$, so
$f(x)\leq c(x,y)-f^c(y)$ and therefore $f\leq f^{c\bar c}$. Applying the
order-reversing transform gives $f^c\geq f^{c\bar c c}$, while the same
majorization applied to $f^c$ gives the reverse inequality. This proves the
triple-transform identity and the fixed-point characterization.

Finally, if $h$ is $c$-concave and $h\geq u$, then
$h^{\bar c}\leq u^{\bar c}$. Transforming again and using
$h=h^{\bar c c}$ gives $h\geq u^{\bar c c}$.
\end{proof}

\subsection{The min-plus cone that is always present}

Ordinary convexity is exceptional, but a nonlinear analogue of conic structure
is automatic.

\begin{proposition}[Universal min-plus structure]\label{prop:tropical}
Let $(u_i)_{i\in I}\subset\Cfrak_c(Y)$ and $(a_i)_{i\in I}\subset\R$. Then,
whenever the right-hand side is proper,
\[
  y\longmapsto\inf_{i\in I}\{u_i(y)+a_i\}
  \quad\text{belongs to }\Cfrak_c(Y).
\]
In particular, $\Cfrak_c(Y)$ is closed under arbitrary pointwise infima and
under addition of constants.
\end{proposition}

\begin{proof}
Write $u_i=f_i^c$. Since $(f_i-a_i)^c=u_i+a_i$, setting
$F(x)=\sup_i\{f_i(x)-a_i\}$ gives directly
\[
  \inf_i(u_i+a_i)
  =\inf_{i,x}\{c(x,\cdot)-f_i(x)+a_i\}
  =\inf_x\{c(x,\cdot)-F(x)\}=F^c.
\]
Thus the lower envelope is a $c$-transform for the original cost on the
original domain $X$.
\end{proof}

\begin{remark}[Tropical, not ordinary, convexity]
In min-plus algebra, ``addition'' is pointwise minimum and scalar
``multiplication'' is addition of a constant. Proposition~\ref{prop:tropical}
says that every $c$-concave class is a tropical convex cone. This statement
does not imply closure under the ordinary average $(u_0+u_1)/2$.
\end{remark}

\begin{example}[A finite cost for which ordinary convexity fails]\label{ex:finite-nonconvex}
Take $X=\{1,2\}$, $Y=\{1,2,3\}$ and
\[
  C=\begin{pmatrix}0&0&2\\0&2&0\end{pmatrix}.
\]
The two rows
$u_0=(0,0,2)$ and $u_1=(0,2,0)$ are $c$-concave: each is obtained by
suppressing the other row with a sufficiently negative source potential. Their
average $\bar u=(0,1,1)$ is not $c$-concave. Indeed, every finite
$c$-concave vector has the form
\[
  \bigl(\min\{A,B\},\min\{A,B+2\},\min\{A+2,B\}\bigr).
\]
If its first coordinate is zero, then $A,B\geq0$ and at least one of them is
zero. If $A=0$, its second coordinate is zero; if $B=0$, its third coordinate
is zero. Neither possibility gives $(0,1,1)$.
\end{example}

The following elementary obstruction explains why ordinary cone structure is
rare for costs on compact spaces.

\begin{proposition}[A uniform modulus obstructs conic scaling]\label{prop:cone-obstruction}
Assume all sections $y\mapsto c(x,y)$ are $L$-Lipschitz, with the same finite
$L$. Then every finite $u\in\Cfrak_c(Y)$ is $L$-Lipschitz. If the class
contains a nonconstant function, it cannot be an ordinary convex cone.
\end{proposition}

\begin{proof}
An infimum of functions sharing the same Lipschitz constant is $L$-Lipschitz.
If $u$ is nonconstant and $\lambda>L/\Lip(u)$, then $\lambda u$ has Lipschitz
constant larger than $L$ and therefore cannot be $c$-concave.
\end{proof}

\section{Three model costs}

The bilinear, quadratic, and metric costs give three distinct answers: a cone,
an affine convex class, and a norm ball.

\subsection{Bilinear cost: the cone of concave functions}

\begin{theorem}[Bilinear cost]\label{thm:bilinear}
Let $X=Y=\R^d$ and $c(x,y)=-\langle x,y\rangle$. Then the proper
upper-semicontinuous $c$-concave functions are exactly the proper
upper-semicontinuous concave functions. The finite-valued part of this class is
an ordinary convex cone, with the constants as a lineality subspace.
\end{theorem}

\begin{proof}
If $u=f^c$, then
\[
  -u(y)=\sup_x\{\langle x,y\rangle+f(x)\}=(-f)^*(y),
\]
so $-u$ is closed convex. Conversely, if $h=-u$ is proper lower-semicontinuous
convex, Fenchel--Moreau gives $h=h^{**}$ \citep{Rockafellar1970}. Choosing
$f=-h^*$ yields
\[
  f^c(y)=\inf_x\{-\langle x,y\rangle+h^*(x)\}=-h^{**}(y)=u(y).
\]
Nonnegative linear combinations of finite concave functions are concave, which
gives the cone statement.
\end{proof}

\begin{remark}[Restricted slope domains]
The full-space hypothesis matters. If $X$ is restricted, the admissible
supporting slopes of $u$ lie in $-X$. When $X$ is convex, the resulting class
is still convex under standard support/closedness assumptions, but it is a cone
only if the allowed slope set is itself stable under nonnegative scaling. Thus
even the bilinear cost need not produce a cone on bounded domains.
\end{remark}

\subsection{Quadratic cost: semiconcavity at a fixed scale}

Let $c_\tau(x,y)=\lvert x-y\rvert^2/(2\tau)$ with $\tau>0$, and set
$q_\tau(y)=\lvert y\rvert^2/(2\tau)$.

\begin{theorem}[Quadratic cost]\label{thm:quadratic}
On $X=Y=\R^d$, a proper upper-semicontinuous function $u$ is
$c_\tau$-concave if and only if $u-q_\tau$ is concave. Equivalently, in the
distributional sense,
\[
  D^2u\preccurlyeq \tau^{-1}I.
\]
The class is an ordinary convex set, but not an ordinary cone.
\end{theorem}

\begin{proof}
Expanding the square gives
\[
  f^{c_\tau}(y)-q_\tau(y)
  =\inf_x\left\{\frac{\lvert x\rvert^2}{2\tau}-f(x)
                 -\frac{\langle x,y\rangle}{\tau}\right\},
\]
an infimum of affine functions of $y$, hence a closed concave function. The
converse follows from Theorem~\ref{thm:bilinear} after rescaling the slope.
The set is convex because $u\mapsto u-q_\tau$ is affine. It is not a cone:
$u=q_\tau$ belongs to the class, whereas $2u-q_\tau=q_\tau$ is not concave.
\end{proof}

This is the first warning against using ``convex cone'' as a generic
description. The curvature scale $\tau^{-1}$ is fixed by the cost. Addition or
large positive scaling changes that scale.

\begin{remark}[Euclidean balls and products of balls]
On a convex Euclidean ball, intrinsic geodesic distance is ordinary Euclidean
distance. The squared-distance cost is therefore the restriction of the
quadratic model above. The $c$-concave class remains convex, with the additional
constraint that supporting slopes come from the prescribed source ball; it is
not a cone. A product of Euclidean balls gives the same statement in the
product Euclidean space. This elementary case should not be confused with the
nontrivial theorem for products of round spheres in Section~\ref{sec:spheres}.
\end{remark}

\subsection{Metric cost: a Lipschitz ball, not a cone}

\begin{theorem}[Metric cost]\label{thm:metric}
Let $(X,d)$ be any metric space, let $L>0$, and set $c(x,y)=Ld(x,y)$. Then
\[
  \Cfrak_c(X)=\{u:X\to\R:\Lip(u)\leq L\}.
\]
For every such $u$, one has $u^c=-u$. This class is closed, convex, balanced,
and stable under pointwise minimum and maximum, but it is not a cone unless it
contains only constants.
\end{theorem}

\begin{proof}
Every section $y\mapsto Ld(x,y)$ is $L$-Lipschitz, so every $c$-transform is
$L$-Lipschitz. Conversely, if $\Lip(u)\leq L$, then
$u(x)\leq u(y)+Ld(x,y)$, hence $Ld(x,y)-u(x)\geq-u(y)$. Taking $x=y$ shows
that the infimum is exactly $u^c(y)=-u(y)$. Applying the same identity to $-u$
gives $u=(-u)^c$. The structural assertions follow from elementary properties
of the Lipschitz seminorm; non-conicity also follows from
Proposition~\ref{prop:cone-obstruction}.
\end{proof}

Modulo additive constants, this is the closed radius-$L$ ball of the normed
space defined by the Lipschitz seminorm. The union over all $L<\infty$ is the
vector space of Lipschitz functions, but changing $L$ changes the cost. In
particular, for the plain geodesic cost $c=d_M$ on a sphere or on any other
metric space, no curvature theorem is needed: the answer is always the
$1$-Lipschitz ball.

\subsection{What uniform semiconcavity does and does not imply}

The quadratic calculation has a useful one-sided generalization.

\begin{proposition}[Semiconcavity inherited from cost sections]\label{prop:section-semiconcavity}
Let $Y\subset\R^d$ be convex. Suppose that for every $x\in X$, the function
\[
  y\longmapsto c(x,y)-\frac{\Lambda}{2}\lvert y\rvert^2
\]
is concave. Then every finite $c$-concave function is $\Lambda$-semiconcave.
In particular, for $c(x,y)=h(x-y)$ it suffices that
$D^2h\preccurlyeq\Lambda I$.
\end{proposition}

\begin{proof}
Subtracting $\Lambda\lvert y\rvert^2/2$ from a $c$-transform leaves an
infimum of concave functions. Pointwise infima preserve concavity because their
hypographs are intersections of convex sets.
\end{proof}

This proposition gives regularity of individual potentials. It does not, by
itself, show that the whole $c$-concave class is closed under ordinary convex
combinations. Quadratic costs are special because the shifted class is exactly
the full concave cone.

\section{Ordinary convexity and nonnegative cross-curvature}

The preceding examples suggest that ordinary convexity is a property of the
cost, not a formal consequence of the transform. In smooth finite dimensions,
the exact criterion is nonnegative cross-curvature.

\subsection{The differential condition}

Assume $X$ and $Y$ are smooth $d$-dimensional domains, $c\in C^4(X\times Y)$,
the mixed Hessian $D^2_{xy}c$ is invertible, and the usual twist and domain
$c$-convexity hypotheses hold. A $c$-segment $s\mapsto x_s$ with base
$\bar y$ is defined by requiring
\[
  D_yc(x_s,\bar y)
  =(1-s)D_yc(x_0,\bar y)+sD_yc(x_1,\bar y).
\]
For compatible $c$-segments $x_s$ and $y_t$, define the cross-curvature by
\begin{equation}\label{eq:cross-curvature}
  \operatorname{cross}_c(\dot x_0,\dot y_0)
  :=-\left.\frac{\partial^4}{\partial s^2\partial t^2}
       c(x_s,y_t)\right|_{s=t=0}.
\end{equation}
The cost has \emph{nonnegative cross-curvature} when this quantity is
nonnegative for all directions. The weak MTW condition A3w imposes the same
sign only on null directions satisfying the corresponding mixed-orthogonality
condition. Nonnegative cross-curvature is therefore stronger than A3w.
For background on these hypotheses and their role in optimal-transport
regularity, see \citet{Villani2009} and \citet{KimMcCann2010}.

\begin{theorem}[Figalli--Kim--McCann]\label{thm:fkm}
Under the smooth twist, nondegeneracy, compactness, and domain $c$-convexity
hypotheses above, the set $\Cfrak_c(Y)$ of $c$-concave functions is an ordinary
convex set if and only if $c$ has nonnegative cross-curvature
\citep{FigalliKimMcCann2011}.
\end{theorem}

\paragraph{Mechanism of the theorem.}
At a point $y_0$, choose supporting cost sections for two potentials $u_0$ and
$u_1$, indexed by $x_0$ and $x_1$. The $c$-segment $x_s$ based at $y_0$
interpolates their cost slopes. Nonnegative cross-curvature says that the
intermediate sliding mountain lies below the same convex combination of the
endpoint mountains. It therefore supports
$u_s=(1-s)u_0+su_1$ at $y_0$. Repeating this at every point proves
$c$-concavity of $u_s$. Conversely, applying convexity of the potential class
to pairs of individual mountains recovers the cross-curvature inequality.

\begin{remark}[A3w is not enough]
A3w, or Loeper's maximum principle, controls the maximum of two endpoint
mountains and is central to continuity of optimal maps. Convexity of the entire
potential class requires the stronger ``below the chord'' inequality, hence
nonnegative cross-curvature in all directions. Regularity of optimal maps and
ordinary convexity of all potentials are related but distinct properties.
\end{remark}

\begin{remark}[Local supports versus global $c$-concavity]
There is a second structural role for MTW-type assumptions. Under twist,
appropriate $c$-convexity of the domains, and A3w, a potential admitting the
correct local $c$-supports can be promoted to a globally $c$-concave
potential. \citet{LoeperTrudinger2021} give a geometric proof under only
twice-differentiable costs. This local-to-global principle concerns recognition
of one potential; Theorem~\ref{thm:fkm} concerns closure of the whole class
under ordinary averages and requires the stronger NNCC condition.
\end{remark}

\subsection{A nonsmooth formulation}

The smooth tensor is not needed to state the underlying geometry. A path,
which need not be continuous, $x_s$ from $x_0$ to $x_1$, based at $\bar y$,
is a variational $c$-segment if
\begin{equation}\label{eq:variational-c-segment}
  c(x_s,\bar y)-c(x_s,y)
  \leq(1-s)\bigl[c(x_0,\bar y)-c(x_0,y)\bigr]
       +s\bigl[c(x_1,\bar y)-c(x_1,y)\bigr]
\end{equation}
for every $y$. A cost space is NNCC if such a path exists for every admissible
triple $(x_0,x_1,\bar y)$. This definition applies to nonsmooth and
infinite-dimensional spaces.

For arbitrary extended costs, the precise convex class is
\[
  \Cfrak_c^\partial(Y)
  :=\{u:\partial^cu(y)\neq\varnothing
       \text{ at every point where }u(y)\text{ is finite}\}.
\]

\begin{theorem}[Synthetic characterization]\label{thm:synthetic}
A cost space is NNCC if and only if $\Cfrak_c^\partial(Y)$ is closed under
ordinary convex combinations. For finite continuous costs on compact spaces,
every finite $c$-concave function has a supporting section, and this reduces to
ordinary convexity of $\Cfrak_c(Y)$
\citep{LegerTodeschiVialard2025}.
\end{theorem}

The qualification by nonempty $c$-subdifferentials matters on noncompact spaces
and for costs taking the value $+\infty$. It should not be silently omitted
when using the modern infinite-dimensional theorem.

\section{Squared geodesic costs on spheres and products}\label{sec:spheres}

The nontrivial geometric example alluded to in the question concerns squared
geodesic distance, not geodesic distance itself.

\begin{theorem}[Round spheres and their products]\label{thm:spheres}
Let
\[
  M=\prod_{k=1}^N\mathbb S^{d_k}_{r_k}
\]
be a finite product of round spheres with its product Riemannian metric, and
let $c(x,y)=d_M(x,y)^2/2$. Then the finite $c$-concave functions on $M$ form an
ordinary convex set. Unless $M$ is a singleton, this set is not an ordinary
cone.
\end{theorem}

\begin{proof}[Proof and attribution]
The round sphere has nonnegative cross-curvature for squared geodesic distance
\citep{KimMcCann2012}. Direct products preserve nonnegative cross-curvature, so
the same holds for $M$. The cut locus prevents a naive global use of the smooth
formula~\eqref{eq:cross-curvature}. For one sphere, \citet{Sei2013} give a
direct global proof of convexity. For products, the product stability of NNCC
combined with the synthetic characterization in Theorem~\ref{thm:synthetic}
gives the stated conclusion. The detailed cut-locus and regularity analysis for
multiple products of spheres is developed by
\citet{FigalliKimMcCann2013}.

Since $M$ is compact, the sections $y\mapsto d_M(x,y)^2/2$ have a common
Lipschitz bound $\operatorname{diam}(M)$. Proposition~\ref{prop:cone-obstruction}
then rules out conic scaling as soon as a nonconstant potential is present.
\end{proof}

\begin{remark}[What is due to Figalli and collaborators]
The general equivalence between convexity of the potential class and
nonnegative cross-curvature is due to Figalli, Kim and McCann. Convexity on a
single sphere was also proved directly by Sei. Kim and McCann established the
cross-curvature and product/submersion mechanisms, while Figalli, Kim and
McCann treated global regularity on arbitrary finite products of round spheres.
Thus ``products of spheres'' is the likely intended reference. Products of
Euclidean balls are already covered by the quadratic Euclidean calculation.
\end{remark}

\section{Other useful convex classes}

Nonnegative cross-curvature is useful partly because it is stable under several
constructions. This produces examples far beyond Euclidean or spherical costs.

\subsection{Bregman divergences}

\begin{proposition}[Bregman costs are NNCC]\label{prop:bregman}
Let $E$ be a Banach space, let $F:E\to\R$ be convex and Gateaux differentiable,
and set
\[
  D_F(x,y)=F(x)-F(y)-DF(y)[x-y].
\]
If $X\subset E$ is convex and $Y\subset E$ is arbitrary, then
$(X\times Y,D_F)$ is NNCC. Consequently, the supported $D_F$-concave class is
convex. For the reversed divergence $D_F(y,x)$, the same conclusion holds if
$DF(X)$ is convex in $E^*$.
\end{proposition}

\begin{proof}
For $x_s=(1-s)x_0+sx_1$, direct expansion gives
\[
  D_F(x_s,\bar y)-D_F(x_s,y)
  =(DF(y)-DF(\bar y))[x_s]+
    F(y)-F(\bar y)+DF(\bar y)[\bar y]-DF(y)[y],
\]
which is affine in $s$. Hence~\eqref{eq:variational-c-segment} holds with
equality. For the reversed divergence, interpolate linearly in the dual
coordinate $DF(x_s)$.
\end{proof}

This includes squared Hilbert distance and explains why relative entropy,
viewed as an infinite-dimensional Bregman divergence, has the same structural
property.

\subsection{Products, quotients, and pullbacks}

\begin{proposition}[Closure mechanisms for NNCC costs]\label{prop:closures}
The following operations preserve nonnegative cross-curvature.
\begin{enumerate}[label=(\roman*)]
  \item If each $(X_k\times Y_k,c_k)$ is NNCC, then the product cost
  $c(x,y)=\sum_k c_k(x_k,y_k)$ is NNCC.
  \item If $c(x,y)=\norm{F(x)-G(y)}_H^2$ for a Hilbert space $H$ and
  $F(X)$ is convex, then $c$ is NNCC.
  \item Suitable cost submersions, including smooth Riemannian submersions for
  squared-distance costs, preserve NNCC.
\end{enumerate}
\end{proposition}

The first statement explains products of spheres. The third generates quotient
geometries, including complex projective spaces from Hopf fibrations and the
Bures--Wasserstein geometry of positive semidefinite matrices from the quotient
$M\mapsto MM^\top$ of a Hilbert space of matrices
\citep{KimMcCann2012,LegerTodeschiVialard2025}.

\subsection{Lifting the cost to probability measures}

For a lower-semicontinuous cost $c$ define its Kantorovich lift
\[
  \mathcal T_c(\alpha,\beta)
  :=\inf_{\pi\in\Pi(\alpha,\beta)}\int c(x,y)\dd\pi(x,y).
\]

\begin{theorem}[Wasserstein lifting of NNCC]\label{thm:lifting}
Under the standard Polish-space, lower-semicontinuity, boundedness, and
finiteness hypotheses, if $(X\times Y,c)$ is NNCC, then
\[
  \bigl(\Pcal(X)\times\Pcal(Y),\mathcal T_c\bigr)
\]
is NNCC. Under a mild condition ensuring finite competitors, the converse also
holds \citep{LegerTodeschiVialard2025}.
\end{theorem}

For $c(x,y)=\lvert x-y\rvert^2$, this says that the cost
$\mathcal T_c= W_2^2$ on the Wasserstein space is itself NNCC. Thus an
appropriate supported class of $W_2^2$-concave \emph{functionals of measures}
is convex. This is a genuinely infinite-dimensional extension of the
quadratic Euclidean statement, and not merely a reformulation of displacement
convexity.

\subsection{A modern catalogue}

The synthetic theory of \citet{LegerTodeschiVialard2025} establishes NNCC for
the following costs, among others. The convexity conclusion in the last column
always refers to the supported class $\Cfrak_c^\partial$; it applies to the full
finite $c$-concave class when compactness and finiteness make supports
automatic.

\begin{center}
\small
\begin{tabularx}{\linewidth}{@{}>{\raggedright\arraybackslash}p{.26\linewidth}
  >{\raggedright\arraybackslash}p{.39\linewidth}X@{}}
\toprule
Space or cost & Representative formula & Structural source \\
\midrule
Hilbert/Bures--Wasserstein &
$\BW^2(A,B)=\tr A+\tr B-2\tr\!\bigl((A^{1/2}BA^{1/2})^{1/2}\bigr)$ &
Hilbert quotient / cost submersion \\
Relative entropy &
$\KL(\alpha\mid\beta)=\int\log(\dd\alpha/\dd\beta)\dd\alpha$ &
Bregman and mixture geometry \\
Hellinger and Fisher--Rao &
$H^2(\alpha,\beta)$ and $\FR^2(\alpha,\beta)$ &
Spherical/Hilbert lifting \\
Selected Hellinger--Kantorovich costs &
Unbalanced transport over an NNCC cone &
Cost submersion; additional hypotheses are needed \\
$2$-Gromov--Wasserstein &
$\GW_2^2$ on gauged metric-measure spaces &
Lifted quadratic structure \\
Wasserstein-on-Wasserstein &
$W_2^2$ on $\Pcal_2(X)$ when the base is NNCC &
Kantorovich lifting \\
\bottomrule
\end{tabularx}
\end{center}

These examples show that convexity of $c$-concave potentials is not restricted
to costs on vector spaces. They also reinforce the need to distinguish a
convex set from a cone: each squared metric carries a fixed curvature or
regularity scale, so positive scaling is generally unavailable.

\section{Discrete and quantitative refinements}

Two further viewpoints are useful in computation and stability theory.

\subsection{Finite costs as tropical polytopes}

For a matrix $C\in\R^{n\times m}$, a $c$-concave vector is
\[
  u_j=\min_{1\leq i\leq n}\{C_{ij}-f_i\}.
\]
Hence the discrete class is the min-plus span of the rows of $C$. It is a
tropical polytope after quotienting by additive constants. The active indices
\[
  \partial^C u(j)=\{i:u_j=C_{ij}-f_i\}
\]
are the discrete $c$-subdifferential. Example~\ref{ex:finite-nonconvex} shows
that a tropical polytope need not be an ordinary convex polytope. This
viewpoint is useful for discrete dual potentials, Laguerre cells, and active-set
algorithms.

\subsection[Strong c-concavity]{Strong $c$-concavity}

Ordinary $c$-concavity only asks that every point admit a supporting cost
section. A quantitative version asks that contact be isolated with a prescribed
growth modulus.

\begin{definition}[Strong $c$-concavity]
A $c$-concave function $u$ is strongly $c$-concave with modulus
$\omega:[0,\infty)\to[0,\infty)$ if, for every $y$, every
$x\in\partial^cu(y)$, and every $z$,
\[
  c(x,z)-u(z)-\bigl(c(x,y)-u(y)\bigr)
  \geq \omega(d_Y(y,z)).
\]
\end{definition}

When $\omega(r)>0$ for $r>0$, a supporting section has a unique contact point.
For the quadratic cost, this is the strong-concavity inequality for the shifted
potential $u-q_\tau$. Under twist and MTW-type hypotheses, differential
criteria for strong $c$-concavity yield quantitative stability of optimal maps
and plans \citep{GallouetMerigotThibert2026}. This is a different strengthening
from ordinary convexity of the \emph{set} of all potentials: it quantifies the
geometry of one potential.

\section{Summary and diagnostic checklist}

The main answers can be compressed as follows.

\begin{center}
\small
\begin{tabularx}{\linewidth}{@{}>{\raggedright\arraybackslash}p{.23\linewidth}
  >{\raggedright\arraybackslash}X
  >{\centering\arraybackslash}p{.11\linewidth}
  >{\centering\arraybackslash}p{.10\linewidth}@{}}
\toprule
Cost / setting & $c$-concave class & Convex? & Cone? \\
\midrule
Arbitrary $c$ & Lower envelopes of shifted sections & Tropical only & Tropical \\
$-\langle x,y\rangle$ on full spaces & Closed concave functions & Yes & Yes \\
$\lvert x-y\rvert^2/(2\tau)$ & $u-\lvert\cdot\rvert^2/(2\tau)$ concave & Yes & No \\
$Ld(x,y)$ & $\Lip(u)\leq L$ & Yes & No \\
$d_M^2/2$ on round spheres/products & Geometric $c$-concave potentials & Yes & No \\
Smooth general $c$ & Convex iff NNCC, under FKM hypotheses & If NNCC & Usually no \\
Finite matrix $C$ & Min-plus row span & Not generally & Tropical \\
\bottomrule
\end{tabularx}
\end{center}

Before asserting convexity or conicity, one should check:
\begin{enumerate}[label=\arabic*.]
  \item Is the cost $d$ or $d^2/2$? The first gives Lipschitz functions on
  every metric space; the second detects geometry.
  \item Are the source and target domains full vector spaces, convex subsets,
  or compact manifolds? Restricted domains constrain supporting slopes.
  \item Is ``convex'' ordinary or min-plus? Min-plus convexity is universal;
  ordinary convexity is not.
  \item Does every potential have a nonempty $c$-subdifferential? On general
  spaces, the NNCC theorem concerns the supported subclass.
  \item Does the cost have NNCC or only A3w? Only NNCC characterizes convexity
  of the whole potential class.
  \item Is a cone claim compatible with the fixed Lipschitz/semiconcavity scale
  imposed by the cost? Usually it is not.
\end{enumerate}

The central organizing principle is therefore:
\begin{center}
\fcolorbox{otline}{otgray}{%
  \begin{minipage}{.88\linewidth}
  \centering\bfseries
  Lower-envelope structure is universal; ordinary convexity is governed by
  NNCC; ordinary conicity requires an additional scaling symmetry.
  \end{minipage}}
\end{center}

\bibliographystyle{plainnat}
\bibliography{references}

\end{document}
